% Compile with: uplatex file.tex && dvipdfmx file.dvi
%           or: latexmk -uplatex file.tex
%\documentclass[uplatex,dvipdfmx]{jsarticle}
\documentclass[dvipdfmx,a4paper,11pt]{jsarticle}
\usepackage{amsmath, amssymb}
\usepackage{tikz}

\title{scenario\_mm1\_queue}
\date{\today}

\begin{document}
\maketitle

到着率 $\lambda$・サービス率 $\mu$ の M/M/1 待ち行列を扱う。利用率 $\rho=\lambda/\mu$（窓口がふさがっている割合）を定義する。

$\rho = \lambda/\mu$

\begin{equation}
  \rho = \lambda/\mu
  \label{eq:E1}
\end{equation}

平均系内人数 $L=\rho/(1-\rho)$ を定義する。

$L = \rho/(1-\rho)$

\begin{equation}
  L = \rho/(1-\rho)
  \label{eq:E2}
\end{equation}

$\rho$ を $\lambda/\mu$ に置き換え、$L$ を $\lambda,\mu$ だけで表す式にする。

E2 の $\rho$ に $\lambda/\mu$ を代入して

\begin{equation}
  L = -\frac{\lambda}{\lambda-\mu}
  \label{eq:E3}
\end{equation}

Little の法則 $L=\lambda W$ を適用するため、両辺を到着率 $\lambda$ で割り、平均滞在時間 $W=L/\lambda$ を得る。

E3 の両辺を $\lambda$ で割って

\begin{equation}
  \frac{L}{\lambda} = -\frac{1}{\lambda-\mu}
  \label{eq:E4}
\end{equation}

\end{document}